\chapter{シミュレーションを用いた運用最適化}
\label{シミュレーションを用いた運用最適化}
本章では，東京23区における救急事案の発生から救急隊の帰署に至るまでの一連の活動を再現するシミュレーション環境の構築と，シミュレーション環境を用いた救急隊配車方法の実験・検討について述べる．シミュレーション環境は，実データに基づく救急需要の時空間分布を再現するとともに，{第\ref{救急活動時間の包括的分析}章}，{第\ref{活動時間延伸の構造的理解}章}% で得られた知見を反映した現実的なモデルとして設計した．
の分析で明らかになった活動時間の特性を反映した現実的なモデルとして設計した．救急隊配車方法については，東京消防庁をはじめ多くの消防組織で一般的に用いられる直近隊運用をベースラインとし，他のヒューリスティック手法やPPO（Proximal Policy Optimization）を用いた深層強化学習手法の間で，平均RT，重症系RT，閾値達成率等の指標から各配車手法の性能を比較した．

\section{用いるデータ}
\subsection{救急搬送データ}
{第\ref{救急活動時間の包括的分析}章}及び{第\ref{活動時間延伸の構造的理解}章}で使用した2019年から2024年までの救急搬送データのうち，2023年1月以降のデータをシミュレーションの事案生成及び検証に使用した．各事案に含まれる項目は\ref{救急搬送データ}で示す表\ref{tab:timeline} と表\ref{tab:data_elements} のとおりである．
なお，本研究で使用する「傷病程度」は、救急搬送データの「収容所見程度」の項目に基づく医療機関収容後の初診時診断に基づくものである。実際の運用では通報段階でこれを完全に把握することは困難であるが、本研究では配車ロジックの理想的な挙動を検証するため、事案発生時点でこの真の傷病程度が判明しているものと仮定する。また，傷病程度と重症度・緊急度は厳密には異なる定義による概念であるが\cite{kinkyu2004}，取り扱えるデータの制約から同質のものとみなす．
% \que{序論で書くかどうか要検討}

\subsection{消防署所データ}
東京23区内の救急隊を保有する177の消防署所の位置情報（緯度・経度）を使用した．各消防署所には救急隊が最大で2台配置されており，シミュレーションでは日中のみ活動するデイタイム救急隊，特定の時期に増隊する非常用救急小隊を除き，2024年10月17日時点で常時運用される192台の救急隊を運用対象とした．

\subsection{医療機関データ}
医療機関については，国土数値情報\cite{medical}から得られる東京23区内に存在する医療機関のうち，救急搬送データから得られる収容先医療機関名と突合し，搬送実績がある230の救急指定病院のみをシミュレーション上での搬送先の対象とする．医療機関の情報としては位置情報のほかに，病床数，救急医療の対応レベル（3次指定病院または2次指定以下病院）が付与されており，病院選択モデルの構築において，これらの属性情報を考慮した確率的選択を実装した．

\subsection{道路ネットワークデータ}
OpenStreetMap\cite{OSM}から取得した道路ネットワークデータを使用し，OSRM（Open Source Routing Machine）\cite{OSRM}を用いて任意の2地点間の移動時間及び移動距離を計算した．
\clearpage

%%%%%%%%%%%%%%%%%%%%%%  5.2   %%%%%%%%%%%%%%%%%%%%%%%%%%%%%%%%%%%
\section{シミュレーション環境の構築}

\subsection{モジュール設計}
シミュレーション環境は，汎用的なEMSシミュレーション環境のオープンソースパイプラインであるOpenEMS\cite{ong2022}を参考に，離散イベントシミュレーション（Discrete Event Simulation：DES）の枠組みに基づき設計した．ただし，本研究では東京都区部の複雑な道路網を考慮し，空間モデリング手法としてOpenEMSで提案された矩形グリッドではなく後述する六角形グリッドシステムを採用するとともに，離散時間イベントの校正方法についても独自の改良を加えることで，より精緻なシミュレーションの実現を図っている．
% DESは，システムの状態変化がイベント発生時のみに起こるという前提で進められるシミュレーション手法であり，EMSのように救急事案の発生や救急隊の状態遷移が離散的に発生するシステムのモデリングに適している．

\subsubsection{空間モデリング}
\label{空間モデリング}
シミュレーションにおける空間表現には，Uber社が開発した六角形階層型空間インデックスシステムであるH3\cite{H3}を採用した．H3は地球表面を階層的な六角形グリッドで分割するシステムであり，従来の矩形グリッドと比較して以下の利点を持つ．
\begin{itemize}
\item \textbf{等距離性}: 六角形の中心から各頂点までの距離が均一であり，隣接セル間の距離計算が容易
\item \textbf{隣接関係の一貫性}：各セルは常に6つの隣接セルを持ち，矩形グリッドにおける4方向と8方向の混在問題を回避
\item \textbf{階層的解像度}：解像度（resolution）を変更することで，同一のインデックス体系で異なる粒度の分析が可能
\end{itemize}
本研究では，保有する救急搬送データの住所情報の粒度を考慮し，解像度9を採用した．この解像度における各六角形の面積は約0.11\si{km^2}，平均辺長は約174mである．東京23区を対象として，救急事案の発生場所，消防署所位置，病院位置をH3インデックスに変換した結果，合計3,120個のグリッドセルが生成された．

グリッド分割の手順は以下のとおりである：
\begin{enumerate}
  \item 東京23区の境界座標内に存在するすべての救急事案，消防署所，病院の位置座標を取得する
  \item 各位置座標をH3ライブラリのgeo\_point\_to\_h3cell関数を用いてH3グリッドに変換する
  \item 重複を除去して一意なH3インデックスの集合を生成し，各インデックスに対して0から始まる連番の行列インデックスを割り当てる
\end{enumerate}
作成された行列インデックスは，後述する移動時間行列及び移動距離行列のアクセスに使用される．

\subsubsection{イベント型の定義}
シミュレーションでは，救急活動の各フェーズに対応する6種類のイベントを定義した．
\begin{enumerate}
\item \textbf{NEW\_CALL}: 新規救急事案の発生
\item \textbf{DISPATCH}：救急隊への出場指令
\item \textbf{ARRIVE\_SCENE}：現場到着
\item \textbf{DEPART\_SCENE}: 現場出発
\item \textbf{ARRIVE\_HOSPITAL}：病院到着
\item \textbf{AMBULANCE\_AVAILABLE}：救急隊の出場可能状態への復帰
\end{enumerate}

また，上記イベントの遷移に対応し，救急隊は以下の6種類で定義されるいずれかの状態を取る．
\begin{enumerate}
\item \textbf{AVAILABLE}: 出場可能（自署にて待機中）
\item \textbf{DISPATCHED}：出場中
\item \textbf{ON\_SCENE}：現場活動中
\item \textbf{TRANSPORTING}: 搬送中
\item \textbf{AT\_HOSPITAL}：病院滞在中
\item \textbf{RETURNING}：帰署中
\end{enumerate}
状態遷移はイベントの処理に伴って発生し，各状態における滞在時間は実データから推定したパラメータに基づいて確率的に生成される．

\subsection{移動時間の計算と校正}
\label{移動時間の計算と校正}
\subsubsection{移動時間の計算}
グリッドセル間の移動は，OpenStreetMapから取得した道路ネットワークデータと，OSRM（Open Source Routing Machine）を用いて計算した．OSRMはOpenStreetMapの地図データを利用して高速に経路を検索するオープンソースのルーティングエンジンであり，任意の2地点間の移動時間及び移動距離を計算できる．

作成した3,120個のグリッドセル間の全組み合わせ（約970万通り）に対して厳密な経路計算を行うことは，計算コストの観点から現実的ではない．そこで，計算時間を短縮するため，全てのグリッドペアから，救急事案のみと救急事案のみのグリッドペアなど，シミュレーションで想定した移動が生じない組合せを除き，計算する必要があるグリッドの選別を行った．また，救急搬送距離の90パーセンタイルが4.9kmであることを踏まえ，5km以内の組合せについては全件をOSRMで計算し，5kmを超えるグリッドペアについては代表サンプルをOSRMで計算したうえで，Haversine距離\cite{sinnott1984}に基づく推定値により補完した．これらの計算結果を各グリッドセル間の移動時間及び移動距離の行列として事前に保存しておくことで，計算効率を向上させた．

\subsubsection{移動時間の校正}
OSRMで計算された移動時間は，一般車両が想定した標準的な走行時間であり，緊急走行が認められる救急車両の特性を反映していない．救急車両は赤色灯とサイレンを使用する緊急走行を行うことで，赤信号の交差点を進行できるほか，法令で定められた停止義務を免除される等，一般車両から優先して道路を通行することが可能であり，一般車両よりも短時間で移動することが可能である．そこで，本研究では実データの移動時間とOSRM計算値の関係を分析し，回帰モデルによる校正を行った．OpenEMSと同様に，実データの移動時間は上位１\%と下位1\%のデータを異常値として取り除いた．校正は救急活動の各フェーズ（出場\_現着，現発\_病着，引揚\_帰署）毎に個別に行い，線形回帰モデルと対数変換校正モデルの2つの校正モデルを作成し，決定係数が高くなるモデルをそれぞれのフェーズに適用し，移動時間行列として保存した．校正後の移動時間行列を用いてシミュレーションを実行することで，実測値との誤差を最小化した現実的な移動時間予測を実現した．

\subsection{サービス時間のモデリング}
移動時間を伴わない活動時間（救急現場での活動時間，病院到着後の引継ぎ時間等）は実データから推定した確率分布に基づいて生成した．OpenEMSやLamら\cite{Lam2016}の研究では救急活動におけるサービス時間は対数正規分布で良く適合することが報告されている．本研究でも，2023年の救急搬送データを用いて各活動区分の時間分布を分析し，対数正規分布のパラメータを推定した．パラメータ推定の際には，総活動時間が傷病程度によって異なることを示した第\ref{救急活動時間の包括的分析}章（\ref{傷病程度別の総活動時間分析TAT}）及び，第\ref{活動時間延伸の構造的理解}章（\ref{パネルデータ分析}）の分析結果，ならびに季節や時間帯による変動を確認した第\ref{救急活動時間の包括的分析}章（\ref{総活動時間の時間的変動}）の分析結果を踏まえ，サービス時間を傷病程度別に最尤推定し，さらに，月別・時間帯別の調整係数を適用する階層的モデルを構築した．時間帯は深夜（0-6時），午前（6-12時），午後（12-18時），夜間（18-24時）の4区分とし，各区分における基本パラメータからの乖離を調整係数として適用する．調整係数は±5\%の閾値を超える場合にのみ適用し，微細な変動を捨象し主要な傾向のみをモデルに反映させる設計とした．

\subsection{病院選択モデルの構築}
実際の救急活動では，傷病者の状態，搬送先病院の受入可否，救急事案発生場所との距離などを考慮して搬送先が決定される．本シミュレーションにおいては上記の厳密な過程を再現することは困難であるため，実データの搬送実績に基づく確率的モデルにより搬送先の病院を決定した．各救急事案に対して，事案発生グリッドセルと各病院の位置関係，病院の属性（救急医療対応レベル，病床数），及び実データにおける搬送頻度を考慮した確率分布から搬送先を選定する．

\subsubsection{重症系事案の病院選択}
重症及び重篤事案では，高度な医療処置が可能な3次指定病院への搬送が優先される．モデルでは，事案発生位置から15㎞圏内に所在する3次指定病院を候補とし，距離による減衰関数
\[
  \mathrm{score} = \exp\left( -\frac{d}{8.0} \right)
\]
に基づいてスコアを算出する．上位3候補からランダムに選択することで，現実の選択における不確実性を反映した．なお，死亡事案については最も近い病院を選択するように設定した．

\subsubsection{軽症系事案の病院選択}
軽症および中等症事案では，2段階の確率的選択モデルを構築した．第1段階として，実績ベースモデルを適用する．時間帯(4時間刻みの6区分)，曜日種別（平日または休日），傷病程度，発生地点をキーとして，過去の搬送実績から病院ごとの選択確率を算出する．このモデルは，特定の時間帯・地域における搬送パターンを再現することを目的としている．

第2段階として，実績データが存在しない組合せに対しては静的フォールバックモデルを適用する．本モデルでは，距離，病床数，過去の受入実績，病院種別と傷病程度の適合性を考慮したスコアリングにより，以下のように各病院の選択確率を算出し正規化を行う．

\[
  P(\mathrm{hospital}) \propto \mathrm{factor}_{\mathrm{dist}} \cdot \mathrm{factor}_{\mathrm{cap}} \cdot \mathrm{factor}_{\mathrm{acc}} \cdot \mathrm{factor}_{\mathrm{genre}}
\]

ここで，$\mathrm{factor}_{\mathrm{dist}}$ は距離による減衰項であり，傷病程度により減衰の強さを変化させた．具体的には，軽症事案では近距離の病院を強く優先し，中等症事案では中程度の距離重視としたほか，傷病程度毎に最大選択距離（軽症：$8\,\mathrm{km}$，中等症：$12\,\mathrm{km}$）を設定し，これを超える病院は候補から除外した．
また，$\mathrm{factor}_{\mathrm{cap}}$ は病床数に基づく重み付け項である．大規模病院ほど受入能力が高いと仮定し，特に中等症において重要視する設定とした．
$\mathrm{factor}_{\mathrm{acc}}$ は過去の受入実績に基づく調整係数であり，特定の病院への搬送が集中することの無いようにした．
$\mathrm{factor}_{\mathrm{genre}}$は病院種別（3次指定病院または2次指定病院以下）の項であり，軽症系事案においては2次指定病院以下が優先的に選ばれるような現実性を反映した．

以上の病院選択モデルにより，距離・病院能力・過去実績・時間帯・傷病程度など多面的な要因を考慮し，実績ベースの部分は，観測データの傾向を忠実に反映し，フォールバック部分はデータの乏しい条件でも妥当な挙動を保証する役割を担うよう設定した．
これにより，シミュレーション上での搬送先決定が，「完全な決め打ち」でも「単純な最近傍ルール」でもない，現実の運用に近い確率的挙動として表現されるようにした．
\clearpage


%%%%%%%%%%%%%%%%%%%%%%  5.3   %%%%%%%%%%%%%%%%%%%%%%%%%%%%%%%%%%%
\section{配車方法の実装\add{➀：ヒューリスティック手法}}
\subsection{直近隊運用（ベースライン）}

% \subsubsection{直近隊運用（ベースライン）}
直近隊運用は，現在東京消防庁をはじめとした多くの消防組織で用いられる標準的な配車方法である．この方法では，救急要請が発生した際に，出場可能な救急隊の中から事案発生場所に最も近い救急隊を選択して配車する．なお，移動距離は都度計算せず，\ref{空間モデリング}で事前計算された移動距離行列から最も移動距離が短い隊を選択する．

具体的なアルゴリズムは以下のとおりである．救急事案$i$が発生した時点で，出場可能な状態（AVAILABLE）にある全ての救急隊の集合を$\mathcal{A}$とする．各救急隊$a \in \mathcal{A}$について，救急隊が配置されたグリッドから事案が発生したグリッドまでの移動距離$d(a, i)$を移動距離行列から取得し，最も短い移動距離を有する救急隊を直近隊として選択する．
\[
  a^* = \mathop{\mathrm{arg\,min}}_{a \in \mathcal{A}} \, d(a, i)
\]

直近隊運用は実装が単純であり，個々の事案に対して最速の対応を保証するという利点がある．一方で，\cite{Carter1972}，\cite{Jagtenberg2017}の先行研究の指摘にあるように，将来の救急需要を考慮しないため，結果として地域全体のカバレッジが低下し，後続の重症事案への対応が遅延する可能性がある．本研究では，この直近隊運用をベースラインとして，提案手法の効果を評価する．


\subsection{傷病度考慮運用}
\label{傷病度考慮運用}
傷病度考慮運用（Severity-Based Strategy）は，傷病度に応じて配車ロジックを切り替えるヒューリスティック手法である．重症系事案には最短のRTで対応できるよう，直近隊運用と同様に事案発生場所に最も近い救急隊を配車し，軽症系事案にはRTとカバレッジのバランスを考慮した救急隊を配車する．

\subsubsection{重症系事案の配車}
重症系事案では，迅速な現場到着が救命率に直結するため，直近隊運用と同様に最寄りの救急隊を選択する．すなわち，傷病度$s$が重症系に該当する場合，
\[
  a^* = \mathop{\mathrm{arg\,min}}_{a \in \mathcal{A}} \, d(a, i)
\]
として配車を行う．

\subsubsection{軽症系事案の配車}
軽症系事案では，許容可能なRTの上限を満たしつつ，将来の重症系事案への対応能力（カバレッジ）を維持することを目指し，以下の2段階のフィルタリングとスコアリングにより救急隊を選択する．

第1段階として，制限時間内に到着可能な救急隊を候補として抽出する．国際的な基準（通報から現着まで15分）を参考に，軽症系事案では13分を制限時間として設定した．

第2段階として，候補となる各救急隊について，応答時間スコアとカバレッジ損失スコアの加重和により総合スコアを算出する．
\begin{equation}
  \mathrm{Score}(a) = w_t \times \mathrm{TimeScore}(a) + w_c \times \mathrm{CoverageLoss}(a)
\end{equation}
ここで，$\mathrm{TimeScore}(a)$はRTを13分で正規化した値，$\mathrm{CoverageLoss}(a)$は当該救急隊が出場した場合のカバレッジ損失率である．重みパラメータは実験を通じて$w_t = 0.2$，$w_c = 0.8$に設定した．

カバレッジ損失の計算では，当該救急隊が所属する消防署が存在するグリッドの周辺から，6分以内または13分以内に到達可能な救急隊数の変化を評価する．具体的には，各サンプルポイントへの最小RTを比較し，カバレッジ率の変化を算出する．この方式により，周辺地域の救急対応能力への影響が大きい救急隊の出場を抑制し，将来の重症系事案への備えを維持する．


\subsection{MEXCLPヒューリスティック運用}
MEXCLP（Maximum Expected Coverage Location Problem）ヒューリスティック運用は，Jagtenbergらの研究\cite{Jagtenberg2017}に基づくカバレッジ最大化手法である．この手法は，救急隊が稼働中（busy）である確率を考慮した期待カバレッジを最大化することを目的とする．

\subsubsection{MEXCLPの基本概念}
従来のカバレッジモデルでは，ある地点がカバーされているとは，その地点に対して閾値時間内に到着可能な救急隊が少なくとも1台存在することを意味した．しかし，この定義では救急隊が他の事案に対応中である可能性を考慮していない．そこで，MEXCLPではbusy確率$q$を導入し，$k$台の救急隊がカバー範囲内にある場合の期待カバレッジを
\[
  E[\mathrm{Coverage}] = 1 - q^k
\]
として計算する．これは，$k$台すべてがbusyである確率$q^k$を1から引いた値であり，少なくとも1台が対応可能である確率を表す．

\subsubsection{動的配車への適用}
MEXCLPの考え方を動的配車問題に適用する．救急要請が発生した際，各候補救急隊について，配車後の残存カバレッジを計算し，残存カバレッジが最大となる救急隊を選択する．

具体的なアルゴリズムは以下のとおりである．まず，閾値時間（13分）以内に事案発生地点に到着可能な救急隊を候補として抽出する．候補が存在する場合，各候補について，その救急隊を除いた残りの救急隊による期待カバレッジを計算する．
\[
  \mathrm{RemainingCoverage}(a) = \sum_{j} d_j \left( 1 - q^{k_j(a)} \right)
\]
ここで，$d_j$は需要点$j$の需要割合，$k_j(a)$は救急隊$a$を除いた場合に需要点$j$をカバーする救急隊数である．残存カバレッジが最大となる救急隊を選択することで，配車後も地域全体のカバレッジを可能な限り維持する．

\subsubsection{計算効率の最適化}
MEXCLPの計算は，全ての需要点に対して残存カバレッジを評価する必要があるため，計算コストが高い．本実装では，需要の累積90\%をカバーする主要グリッドのみを計算対象とすることで，計算効率を改善した．具体的には，過去の救急事案データから各H3グリッドの需要割合を算出し，需要の高い順にソートして累積需要が90\%に達するまでのグリッドを選択した．これにより，計算対象のグリッド数を約3,120個から数百個に削減しつつ，カバレッジ評価の精度を維持した．また，閾値時間内に到着可能な候補が存在しない場合は，直近隊運用にフォールバックし，最寄りの救急隊を選択する．


%%%%%%%%%%%%%%%%%   5.4       %%%%%%%%%%%%%%%%%%%
\section{\add{配車方法の実装②：}深層強化学習手法}
前節までのヒューリスティック手法では，人が設計したルールに基づいて配車を行う\del{.これに対し，深層強化学習を用いて，データから最適な方策を学習するアプローチを提案する．}\add{ため，ルール設計者の知識や経験に依存し，複雑な状況への適応に限界があると考えられる．本節では，本研究の主要な提案手法として，深層強化学習を用いた配車最適化アプローチを提案する．深層強化学習は，シミュレーション環境との相互作用を通じて，データから最適な配車方策を学習することが可能であり，救急需要の時空間的な変動や救急隊の稼働状況など，多様な要因を考慮した動的な配車判断の実現が期待できる．}

\subsection{マルコフ決定過程としての定式化}
救急車配車問題をMDP（Markov Decision Process）として定式化する．MDPは，状態空間$S$，行動空間$A$，状態遷移確率$P$，報酬関数$R$，割引率$\gamma$の5つの要素$(S, A, P, R, \gamma)$で定義される．

\subsubsection{状態空間}
エージェントが観測する状態$s_t$は，以下の4つの情報から構成される．
\begin{equation}
s_t = ( A_t, I_t, T_t, G_t )
\end{equation}

\paragraph{救急隊情報$A_t$（960次元）}

各救急隊$i$について，表\ref{tab:feature_normalization}に示す5つの特徴量を定義する．状態空間における次元数は$192 \times 5 = 960$次元である．

\begin{table}[H]
\centering
\caption{救急隊情報の特徴量と正規化方法}
\label{tab:feature_normalization}
\begin{tabularx}{\linewidth}{cllX}
\hline
次元 & 特徴量 & 正規化式 & 説明 \\
\hline
1 & 緯度 & $(\mathrm{lat} + 90) / 180$ & 現在位置（南北方向）を$[0, 1]$に正規化 \\
2 & 経度 & $(\mathrm{lng} + 180) / 360$ & 現在位置（東西方向）を$[0, 1]$に正規化 \\
3 & 利用可能性 & $\{0, 1\}$ & 出場可能な場合は1，それ以外は0 \\
4 & 出動回数 & $\min(\mathrm{calls} / 10, 1)$ & 当日の累積出動回数（10回で飽和） \\
5 & 移動時間 & $\min(\mathrm{time} / 1800, 1)$ & 事案現場までの移動時間（30分で正規化） \\
\hline
\end{tabularx}
\end{table}

移動時間は\ref{移動時間の計算と校正}で構築した校正済み移動時間行列を用いて，救急隊の現在位置から事案発生場所までの所要時間を取得し，30分を上限として正規化した．

\paragraph{事案情報$I_t$（10次元）}

現在発生している救急事案の情報を表\ref{tab:incident_feature_normalization}のように定義する．

\begin{table}[H]
\centering
\caption{事案情報の特徴量と正規化方法}
\label{tab:incident_feature_normalization}
\begin{tabularx}{\linewidth}{cllX}
\hline
次元 & 特徴量 & 正規化方法 & 説明 \\
\hline
1--2 & 事案位置 & $[0, 1]$正規化 & 事案発生地点の地理的位置 \\
3--8 & 傷病程度 & one-hot & 傷病程度をバイナリベクトルで表現 \\
9 & 待機時間 & $\min(\mathrm{wait} / 600, 1)$ & 事案発生からの経過時間（10分で飽和） \\
10 & 優先度スコア & $[0, 1]$ & 時間と傷病程度から算出する総合的な優先度 \\
\hline
\end{tabularx}
\end{table}


\paragraph{時間情報$T_t$（8次元）}

救急需要の時間的パターンを表\ref{tab:time_features}のように定義する．

\begin{table}[H]
\centering
\caption{時間情報の特徴量}
\label{tab:time_features}
\begin{tabularx}{\linewidth}{clX}
\hline
次元 & 特徴量 & 説明 \\
\hline
1 & エピソード進行度 & 1日の中での時間経過割合 \\
2--3 & 時刻（周期） & $\sin(2\pi h / 24)$，$\cos(2\pi h / 24)$による24時間周期エンコーディング \\
4--7 & 時間帯フラグ & 朝(6--10時)，昼(10--17時)，夕(17--21時)，夜(21--6時)のone-hot \\
8 & 平日フラグ & 平日の場合は1，休日の場合は0 \\
\hline
\end{tabularx}
\end{table}


\paragraph{大域的情報$G_t$（21次元）}

システム全体の状況を把握するため，表\ref{tab:global_features}に示す統計量を状態に含める．

\begin{table}[H]
\centering
\caption{大域的情報の特徴量}
\label{tab:global_features}
\begin{tabularx}{\linewidth}{clX}
\hline
次元 & 特徴量 & 説明 \\
\hline
1--8 & 利用可能救急車の移動時間統計 & 最短，平均，中央値，最長，標準偏差，第1四分位，第3四分位，分散 \\
9--14 & 全救急車の移動時間統計 & 利用不可の救急車も含めた統計 \\
15--19 & 稼働状況 & 利用可能率，稼働率，絶対数，閾値達成数，利用可能フラグ \\
20--21 & カバレッジ率 & 6分，13分以内に到達可能な地域の割合 \\
\hline
\end{tabularx}
\end{table}

これらを結合した状態ベクトルは999次元となる．


\subsubsection{行動空間}
エージェントの行動空間は，発生した救急事案に対し，時刻$t$における利用可能な救急隊集合の中から1台を選択することとして定義される．
\[
a_t \in \mathcal{A}_t = \{i : \text{ambulance}_i \text{ is available at time } t\}
\]
ここで，$|\mathcal{A}_t|$は時刻$t$における利用可能な救急隊数である．ただし，利用可能でない救急隊への配車は無効な行動となるため，行動マスク$m_t[i]$を用いて有効な行動のみを選択可能とする．具体的には，各時刻$t$において，行動マスクを以下のように定義する．
\[
  m_t[i] =
  \begin{cases}
    1 & \text{if } i \in \mathcal{A}_t \\
    0 & \text{otherwise}
  \end{cases}
\]


\subsubsection{報酬関数}
報酬関数は，RTの短縮を促進するよう設計した．本研究では，傷病度（重症系・軽症系）に応じた目標時間に基づく連続的な報酬関数を設計した．

\paragraph{重症系事案の報酬関数}

重症系事案では，6分以内の到着が期待されるため，指数関数的に減衰する報酬を設計した．
\[
r_{\text{critical}}(t) = 
\begin{cases} 
B_c \cdot \exp(-\lambda t) & \text{if } t \leq T_c \\ 
-\alpha_c \cdot (t - T_c)^{1.5} & \text{if } t > T_c 
\end{cases}
\]
ここで，$t$はRT（分），$T_c(=6)$は目標時間，$B_c$は最大ボーナス，$\lambda$は減衰パラメータ，$\alpha_c$はペナルティスケールである．指数関数的な報酬により，1分でも早い到着に対して大きな報酬を与える設計とした．また，目標時間超過時には1.5乗のペナルティを課すことで，大幅な遅延に対するペナルティが増加するようにした．

\paragraph{軽症系事案の報酬関数}

軽症系事案では，13分以内の到着を目標とし，線形的な報酬を適用した．
\[
r_{\text{mild}}(t) = 
\begin{cases} 
B_m \cdot (1 - t/T_m) & \text{if } t \leq T_m \\ 
-\alpha_m \cdot (t - T_m) & \text{if } t > T_m 
\end{cases}
\]
ここで，$t$はRT（分），$T_m(=13)$は目標時間，$B_m$は最大ボーナス，$\alpha_m$はペナルティスケールである．軽症系事案のペナルティは重症系事案と比較して小さく設定することで，RTがやや長くなることを許容しつつ，システム全体のカバレッジ維持を促進する設計とした．


\subsection{PPOアルゴリズム}
方策の最適化には，Proximal Policy Optimization（PPO）アルゴリズム\cite{schulman2017}\cite{Moen2024}を採用した．PPOは方策勾配法の一種であり，方策の更新幅を制限することで学習の安定性を確保する手法である．

PPOでは，救急事案発生時刻$t$における最適な配車決定$a_t$を求めるため，以下の目的関数を最大化する．
\[
J(\theta) = \mathbb{E}_{\tau \sim \pi_\theta} \left[ \sum_{t=0}^{T} \gamma^t r(s_t, a_t) \right]
\]
ここで，$\theta$は方策パラメータ，$\tau$は軌跡，$\gamma$は割引率，$r(s_t, a_t)$は報酬関数である．

方策パラメータの更新には，以下のPPOクリップ目的関数を用いる．
\[
L^{\text{CLIP}}(\theta) = \hat{\mathbb{E}}_t \left[ \min\left( \rho_t(\theta) \hat{A}_t, \text{clip}(\rho_t(\theta), 1-\epsilon, 1+\epsilon) \hat{A}_t \right) \right]
\]
ここで，$\rho_t(\theta) = \frac{\pi_\theta(a_t|s_t)}{\pi_{\theta_{\text{old}}}(a_t|s_t)}$は方策比率，$\hat{A}_t$はアドバンテージ推定値，$\epsilon$はクリッピング係数である．

\subsubsection{ネットワーク構造}
方策ネットワーク（Actor）及び価値ネットワーク（Critic）は，表\ref{tab:network_structure}に示す構造を持つ多層パーセプトロンとして実装した．

\begin{table}[htbp]
\centering
\caption{ネットワーク構造}
\label{tab:network_structure}
\begin{tabular}{lll}
\hline
項目 & Actor & Critic \\
\hline
入力層 & 状態次元（999または37） & 状態次元 \\
隠れ層1 & 256ユニット + ReLU + Dropout(0.1) & 256ユニット + ReLU + Dropout(0.1) \\
隠れ層2 & 128ユニット + ReLU + Dropout(0.1) & 128ユニット + ReLU + Dropout(0.1) \\
出力層 & 行動次元（ソフトマックス） & 1（状態価値） \\
初期化 & Xavier均一初期化 & スケール0.001の正規分布 \\
\hline
\end{tabular}
\end{table}
\clearpage

\subsubsection{ハイパーパラメータ}
主要なハイパーパラメータを表\ref{tab:ppo_hyperparameters}に示す．

\begin{table}[htbp]
\centering
\caption{PPOの主要なハイパーパラメータ}
\label{tab:ppo_hyperparameters}
\begin{tabular}{lll}
\hline
パラメータ & 値 & 説明 \\
\hline
学習率（Actor） & 0.0003--0.003 & 方策ネットワークの学習率 \\
学習率（Critic） & 0.001--0.01 & 価値ネットワークの学習率 \\
バッチサイズ & 1024--2048 & ミニバッチサイズ \\
エポック数 & 6--10 & 各バッチの学習回数 \\
割引率$\gamma$ & 0.99 & 将来報酬の割引率 \\
GAE $\lambda$ & 0.95 & Generalized Advantage Estimationパラメータ \\
エントロピー係数 & 0.015 & 探索促進のためのエントロピーボーナス \\
勾配クリッピング & 0.5 & 勾配の最大ノルム \\
\hline
\end{tabular}
\end{table}
\clearpage


\subsection{提案手法のバリエーション}
本研究では，配車問題の特性に応じた3つのPPOバリエーションを提案・比較する．

\subsubsection{（1）通常PPO（Full State Space）}
前述の999次元状態空間と192次元行動空間を持つ標準的なPPOモデルである．すべての救急隊の情報を考慮し，システム全体の最適化を目指す．

\subsubsection{（2）ハイブリッドPPO}
本研究における重症系事案に対する最適解は，個別事案においては最短のRTが保証される直近隊運用であることが先行研究\cite{Carter1972}からも示されている．そこで，重症系事案への最短対応と軽症系事案の効率的な配車を両立するため，傷病度に応じて配車ロジックを切り替えるハイブリッド戦略を提案する．

戦略：
\begin{itemize}
  \item 重症系事案（重症・重篤・死亡）：直近隊運用を適用し，最短のRTを保証
  \item 軽症系事案（軽症・中等症）：PPOによる学習済み方策を適用し，カバレッジを考慮した配車
\end{itemize}

\paragraph{報酬関数}

ハイブリッドモードでは，軽症系事案に対してのみ報酬を計算し，RTとカバレッジをそれぞれ重み付けして統合する．
\[
r_{\text{hybrid}} = w_{RT} \cdot r_{\text{time}} + w_{cov} \cdot r_{\text{coverage}}
\]
\begin{itemize}
  \item RT報酬$r_{\text{time}}$：1分経過ごとにペナルティ，13分以内達成でボーナス
  \item カバレッジ報酬$r_{\text{coverage}}$：カバレッジ率0.8以上で大きいボーナス，0.6--0.8で小さいボーナス，0.6未満でペナルティ
\end{itemize}

\subsubsection{（3）コンパクトPPO（Compact State Space）}
大規模な状態空間による学習の困難さを軽減するため，状態空間を大幅に削減したコンパクトモデルを提案する．

\paragraph{Top-K状態表現}

全192台の救急隊情報を保持する代わりに，事案発生地点に対して移動時間が短い上位$K$台（$K=10$）のみを考慮する．
\[
A_t^{\text{compact}} = \underset{i \in \mathcal{A}_t}{\operatorname{Top-K}} \left( \text{travel\_time}(i, \text{incident}) \right)
\]

\paragraph{状態空間（$K=10$の場合，37次元）}

\begin{table}[H]
\centering
\caption{コンパクトPPOの状態空間（$K = 10$の場合）}
\label{tab:compact_state_space}
\begin{tabular}{cll}
\hline
カテゴリ & 次元数 & 内容 \\
\hline
救急隊情報 & $K \times 3 = 30$ & 移動時間，移動距離，カバレッジ損失 \\
事案情報 & 2 & 傷病度（軽症系/重症系のバイナリ） \\
時間特徴 & 1 & 時刻の周期エンコーディング \\
大域的状態 & 4 & 利用可能台数，カバレッジ率，6分圏内台数，平均移動時間 \\
\hline
\end{tabular}
\end{table}

\paragraph{行動空間（$K$次元）}

Top-K救急隊の中から1台を選択する．

この削減により，状態次元を999から37に圧縮し，学習効率の向上を図る．
\clearpage
%%%%%%%%%%%%%%%%%%%%%%%%%  5.4   %%%%%%%%%%%%%%%%%%%%%%%%%%%%%%%%%%%%%%%%%
\section{実験設定}
本節では，提案する配車方法の評価実験の設定について述べる．

\subsection{実験対象期間}
配車方法の性能を多様な条件下で検証するため，2024年から表\ref{tab:シミュレーション対象期間}に示す6つの代表的な週を選定した．選定にあたっては，図\ref{fig:件数と重症割合}に示すとおり事案件数と重症率の2軸により4象限に分類し，各象限から少なくとも一つの期間を含めることで，季節性や需要特性の異なる条件下での性能を評価できるようにした．

\begin{table}[htbp]
    \centering
    \caption{シミュレーション対象期間の選定理由と特性}
    \label{tab:シミュレーション対象期間}
    \begin{tabular}{clll}
        \toprule
        テスト週 & 季節 & 特性 & 選定理由 \\
        \midrule
        2024-02-04～02-10 & 冬 & 低件数$\times$高重症率 & 冬季の閑散期（重症割合最高） \\
        2024-03-31～04-06 & 春 & 低件数$\times$高重症率 & 閑散期内の高重症率 \\
        2024-05-05～05-11 & 春 & 低件数$\times$低重症率 & 閑散期内の低重症率 \\
        2024-06-30～07-06 & 夏 & 高件数$\times$低重症率 & 夏季の繁忙期開始（重症割合最低） \\
        2024-07-21～07-27 & 夏 & 高件数$\times$低重症率 & 夏季の最繁忙期 \\
        2024-12-22～12-28 & 冬 & 高件数$\times$高重症率 & 冬季の最繁忙期（搬送件数最大） \\
        \bottomrule
    \end{tabular}
\end{table}

\begin{figure}[H]
\centering
\includegraphics[width=\textwidth]{chap5-図/2024の件数と重症割合-週次-脚注.png}
\caption{2024年の週次搬送件数と重症割合}
\label{fig:件数と重症割合}
\end{figure}

各テスト期間において，「当該週の中からランダムに1日を選択し，24時間のシミュレーションを実行する」という試行を各配車手法について100回繰り返し，評価指標の平均値を比較した．

\subsection{学習設定}
PPOモデルの学習には2023年のデータを使用した．学習期間と評価期間については多様な設定で実験を行い，各モデルの中で比較的好成績だった設定を表\ref{tab:モデルの学習期間と設定}に示す．エピソード数は学習の収束状況から，適切な学習回数を設定した．

\begin{table}[H]
    \centering
    \caption{各モデルの学習期間および設定}
    \label{tab:モデルの学習期間と設定}
    \resizebox{\columnwidth}{!}{%
        \begin{tabular}{llcll}
            \toprule
            モデル & 学習期間 & 評価期間 & エピソード数 & 備考 \\
            \midrule
            通常PPO & 2023/1/15--2/4 & 2023/2/5--2/11 & 8,000 & 冬季の重症系事案への対応を重視 \\
            ハイブリッドPPO & 2023/1/15--2/4 & 2023/2/5--2/11 & 8,000 & 冬季の重症系事案への対応を重視 \\
            コンパクトPPO & \makecell[l]{2023/4/9--5/14，7/9--30，\\11/19--12/17} & 2023/9/17--10/15 & 3,000 & 多様な季節を学習 \\
            \bottomrule
        \end{tabular}%
    }
\end{table}
1エピソードは1日（24時間）のシミュレーションに相当し，その日に発生する全ての救急事案に対する配車決定を学習する．

\subsection{計算環境}
実験は以下の計算環境で実施した：
\begin{itemize}
  \item GPU：NVIDIA GeForce RTX 4070 Ti SUPER
  \item CPU：Intel Core i7-14700F
  \item メモリ：64GB DDR5-5600MT/s
  \item 実行環境：Python3.12，PyTorch2.7
\end{itemize}

\subsection{比較手法}
以下の配車手法を比較対象とした：

\noindent\textbf{ヒューリスティック手法}
\begin{enumerate}[label=\arabic*., leftmargin=*]
    \item 直近隊運用：移動距離行列に基づいた最短距離で到着可能な救急隊を選択
    \item 傷病度考慮運用：\ref{傷病度考慮運用}で述べた傷病程度に応じた配車戦略
    \item MEXCLPヒューリスティック：既存研究\cite{Jagtenberg2017}で提案されたカバレッジを考慮した戦略（以下，MEXCLPと記載する）
\end{enumerate}

\noindent\textbf{深層強化学習手法}
\begin{enumerate}[resume, label=\arabic*., leftmargin=*]
    \item 通常PPO：5.3.2節で定義した999次元の状態空間と192次元の行動空間を持つ標準的なPPOモデル．全救急隊の情報を考慮し，重症系・軽傷系の区別なく統一的な方策を学習する．
    \item ハイブリッドPPO：重症系事案には直近隊運用を適用し，軽傷系事案にのみPPOによる方策を適用するハイブリッド手法．軽傷系事案においてカバレッジ維持を重視した学習を行う．
    \item コンパクトPPO：状態空間をTop-Kの救急隊情報に限定し，次元を圧縮したモデル．学習効率の向上を図りつつ，実用的な配車性能を目指す．
\end{enumerate}

\subsection{評価指標}
各配車手法の性能を以下の指標で評価した：

\begin{table}[htbp]
    \centering
    \caption{評価指標の定義と目標}
    \label{tab:evaluation_metrics}
    \begin{tabular}{llc}
        \toprule
        指標 & 定義 & 目標 \\
        \midrule
        重症系平均RT & 重症系事案（重症・重篤・死亡）の平均RT（分） & 最小化 \\
        全体平均RT & 全事案の平均RT（分） & 最小化 \\
        % 6分達成率（全体） & 全事案のうち6分以内に現着した割合 (\%) & 最大化 \\
        6分達成率（重症系） & 重症系事案のうち6分以内に現着した割合 (\%) & 最大化 \\
        13分達成率（全体） & 全事案のうち13分以内に現着した割合 (\%) & 最大化 \\
        \bottomrule
    \end{tabular}
\end{table}
\clearpage
%%%%%%%%%%%%%%%%%%%%%%%%%  5.5   %%%%%%%%%%%%%%%%%%%%%%%%%%%%%%%%%%%%%%%%%
\section{実験結果}
本節では，各配車手法の比較実験結果を示す．

\subsection{重症系事案のレスポンスタイム}
表\ref{tab:重症系事案のRT}に各テスト期間における重症系事案の平均RTを示す．

重症系事案のRTについては，傷病度考慮運用が低件数期間（2月，3月，5月）において最良の結果を示し，直近隊運用と比較して0.5～0.7分の改善が確認された．これは，軽傷系事案に対してカバレッジを考慮した配車を行うことで，重症系事案が発生した際により近い救急隊が待機している状態を維持できたためと考えられる．

一方，高件数期間（6月，7月，12月）では傷病度考慮運用の優位性が低下し，6月ではコンパクトPPOが8.62分で最長となった．繁忙期においては，救急隊の稼働率が高く，カバレッジを考慮した配車を行う余裕が限られるため，直近の救急隊を配車する単純な戦略が有効であることが示唆される．

コンパクトPPOは全期間を通じて直近隊運用と近い値を示した．ハイブリッドPPOは重症系事案に直近隊運用を適用する設計であるにもかかわらず，直近隊運用より0.1～0.8分遅れる結果となった．これは，軽症系事案の対応遅延が間接的に重症系事案のRTに影響を与えた可能性を示唆している．

通常PPO及びMEXCLPは全期間で最も遅いRTを示し，重症系事案への対応には適さないことが確認された．

\begin{sidewaystable}[p]
    \centering
    \caption{各テスト期間における配車手法別の重症RT（分）}
    \label{tab:重症系事案のRT}
    \resizebox{\textwidth}{!}{
    \begin{tabular}{lcccccc}
        \toprule
        テスト週 & 2024/2/4 & 2024/3/31 & 2024/5/5 & 2024/6/30 & 2024/7/21 & 2024/12/22 \\
        \makecell[l]{季節\\象限特性} & \makecell{冬\\低件数$\times$高重症} & \makecell{春\\低件数$\times$高重症} & \makecell{春\\低件数$\times$低重症} & \makecell{夏\\高件数$\times$低重症} & \makecell{夏\\高件数$\times$低重症} & \makecell{冬\\高件数$\times$高重症} \\
        \midrule
        直近隊運用 & 8.03 & 7.37 & 7.35 & 8.73 & 10.61 & \textbf{10.88} \\
        傷病度考慮 & \textbf{7.38} & \textbf{6.76} & \textbf{6.81} & 9.06 & \textbf{10.34} & 10.95 \\
        MEXCLP & 10.74 & 10.66 & 10.54 & 11.71 & 12.60 & 12.61 \\
        通常PPO & 10.02 & 9.13 & 9.50 & 11.45 & 13.14 & 14.01 \\
        ハイブリッドPPO & 8.23 & 7.55 & 7.53 & 9.57 & 10.74 & 11.29 \\
        コンパクトPPO & 7.89 & 7.21 & 7.36 & \textbf{8.62} & 10.61 & \textbf{10.88} \\
        \bottomrule
    \end{tabular}
    }
    \footnotesize
    \begin{flushleft}
    注：各値は重症系事案における平均RT（分）を示す．
    \end{flushleft}
\end{sidewaystable}
\clearpage

\subsection{全体のレスポンスタイム}
表\ref{tab:全体のRT}に全事案の平均RTを示す．

全体のRTについては，コンパクトPPOが6期間中5期間で最良の結果を示した．直近隊運用と比較して，閑散期では0.07～0.19分，繁忙期（6月）では0.16分の改善が確認された．唯一7月期のみ直近隊運用が最良（10.81分）となったが，コンパクトPPOとの差は0.04分と僅差であった．

傷病度考慮運用は重症系RTでは優れた結果を示す一方，全体RTでは直近隊運用と比較して2.0～2.5分の大幅な悪化が確認された．これは，軽症系事案（全事案の約9割）に対してカバレッジを優先した遠方の救急隊を配車するためである．

ハイブリッドPPOは全手法中で最も遅い全体RTを示した．軽傷系事案に対するカバレッジ重視の学習が過度に働き，不必要に遠方の救急隊を選択する傾向が生じたと考えられる．通常PPOも同様に，999次元の状態空間からの学習が望ましい方向に収束せず，効率的な配車方策を獲得できなかったことが可能性がある．

\begin{sidewaystable}[p]
    \centering
    \caption{各テスト期間における配車手法別の全体RT（分）}
    \label{tab:全体のRT}
    \resizebox{\textwidth}{!}{
    \begin{tabular}{lcccccc}
        \toprule
        テスト週 & 2024/2/4 & 2024/3/31 & 2024/5/5 & 2024/6/30 & 2024/7/21 & 2024/12/22 \\
        \makecell[l]{季節\\象限特性} & \makecell{冬\\低件数$\times$高重症} & \makecell{春\\低件数$\times$高重症} & \makecell{春\\低件数$\times$低重症} & \makecell{夏\\高件数$\times$低重症} & \makecell{夏\\高件数$\times$低重症} & \makecell{冬\\高件数$\times$高重症} \\
        \midrule
        直近隊運用 & 7.94 & 7.42 & 7.51 & 9.46 & \textbf{10.81} & 11.32 \\
        傷病度考慮 & 9.95 & 9.82 & 9.93 & 11.57 & 12.50 & 12.96 \\
        MEXCLP & 10.71 & 10.60 & 10.54 & 11.92 & 12.56 & 12.75 \\
        通常PPO & 9.88 & 9.04 & 9.27 & 11.37 & 13.08 & 13.74 \\
        ハイブリッドPPO & 11.08 & 10.40 & 10.51 & 12.68 & 13.78 & 14.05 \\
        コンパクトPPO & \textbf{7.87} & \textbf{7.25} & \textbf{7.32} & \textbf{9.30} & 10.85 & \textbf{11.24} \\
        \bottomrule
    \end{tabular}
    }
    \footnotesize
    \begin{flushleft}
    注：各値は全事案における平均RT（分）を示す．
    \end{flushleft}
\end{sidewaystable}
\clearpage

\subsection{6分閾値達成率（重症系）}
表\ref{tab:6分閾値達成率}に6分閾値達成率を示す．

6分閾値達成率については，傷病度考慮運用が全期間で最良の結果を示し，直近隊運用と比較して3.4～8.9ポイントの改善を達成した．特に閑散期（3月，5月）では44.6\%と最も高い達成率を記録した．

コンパクトPPOは直近隊運用とほぼ同等の達成率を示し，一部の期間（2月，6月，12月）では直近隊運用を1.3～1.7ポイント上回った．

MEXCLPは全期間で3～5\%台と極めて低い達成率となった．これは，MEXCLPがカバレッジ最大化を目的としており，個々の事案への迅速な対応よりも地域全体の期待カバレッジを優先するためである．

通常PPOは10～18\%程度の達成率にとどまり，重症系事案への迅速対応を十分に学習できていないことが確認された．

\begin{sidewaystable}[p]
    \centering
    \caption{重症系事案の6分閾値達成率（\%）}
    \label{tab:6分閾値達成率}
    \resizebox{\textwidth}{!}{
    \begin{tabular}{lcccccc}
        \toprule
        テスト週 & 2024/2/4 & 2024/3/31 & 2024/5/5 & 2024/6/30 & 2024/7/21 & 2024/12/22 \\
        \makecell[l]{季節\\象限特性} & \makecell{冬\\低件数$\times$高重症} & \makecell{春\\低件数$\times$高重症} & \makecell{春\\低件数$\times$低重症} & \makecell{夏\\高件数$\times$低重症} & \makecell{夏\\高件数$\times$低重症} & \makecell{冬\\高件数$\times$高重症} \\
        \midrule
        直近隊運用 & 31.6 & 37.7 & 35.7 & 27.7 & 22.5 & 21.9 \\
        傷病度考慮 & \textbf{37.0} & \textbf{44.6} & \textbf{44.6} & \textbf{31.5} & \textbf{23.7} & \textbf{25.5} \\
        MEXCLP & 3.8 & 3.6 & 5.1 & 4.5 & 4.2 & 3.2 \\
        通常PPO & 15.0 & 17.9 & 18.0 & 12.2 & 10.0 & 10.9 \\
        ハイブリッドPPO & 29.5 & 38.1 & 34.5 & 24.3 & 21.6 & 20.2 \\
        コンパクトPPO & 33.0 & 35.9 & 34.5 & 29.0 & 22.4 & 23.6 \\
        \bottomrule
    \end{tabular}
    }
    \footnotesize
    \begin{flushleft}
    注：各値は重症系事案における6分以内到着率（\%）を示す．
    \end{flushleft}
\end{sidewaystable}
\clearpage

\subsection{13分閾値達成率（全体）}
表\ref{tab:13分閾値達成率}に13分閾値達成率を示す．

13分閾値達成率については，MEXCLPが6期間中5期間で最良の結果を示し，特に5月期では97.1\%を達成した．これは既存研究\cite{Jagtenberg2017}の知見と一致し，MEXCLPが平均RTを犠牲にしつつも閾値以内到着率を最大化する特性を持つことが確認できた．

直近隊運用とコンパクトPPOは僅差の達成率を示し，両手法とも閑散期で90\%以上，繁忙期では67～81\%を達成した．両手法の差は0.4ポイント以内であり，実質的に同等の性能と評価できる．

傷病度考慮運用とハイブリッドPPOは全期間で大幅に低い達成率を示した．特に繁忙期（12月）では54\%台まで低下し，全事案の約半数が13分以内に到着できない状況となった．

\begin{sidewaystable}[p]
    \centering
    \caption{全体の13分閾値達成率（\%）}
    \label{tab:13分閾値達成率}
    \resizebox{\textwidth}{!}{
    \begin{tabular}{lcccccc}
        \toprule
        テスト週 & 2024/2/4 & 2024/3/31 & 2024/5/5 & 2024/6/30 & 2024/7/21 & 2024/12/22 \\
        \makecell[l]{季節\\象限特性} & \makecell{冬\\低件数$\times$高重症} & \makecell{春\\低件数$\times$高重症} & \makecell{春\\低件数$\times$低重症} & \makecell{夏\\高件数$\times$低重症} & \makecell{夏\\高件数$\times$低重症} & \makecell{冬\\高件数$\times$高重症} \\
        \midrule
        直近隊運用 & 91.5 & 94.7 & 94.3 & 80.5 & 70.9 & 67.4 \\
        傷病度考慮 & 81.4 & 76.6 & 75.6 & 63.7 & 57.4 & 54.1 \\
        MEXCLP & \textbf{94.3} & \textbf{95.8} & \textbf{97.1} & 80.4 & \textbf{73.0} & \textbf{70.9} \\
        通常PPO & 82.9 & 88.7 & 87.0 & 72.8 & 61.4 & 56.9 \\
        ハイブリッドPPO & 74.7 & 80.2 & 79.2 & 64.2 & 56.6 & 54.6 \\
        コンパクトPPO & 91.1 & 95.3 & 94.7 & \textbf{81.1} & 70.3 & 67.5 \\
        \bottomrule
    \end{tabular}
    }
    \footnotesize
    \begin{flushleft}
    注：各値は全事案における13分以内到着率（\%）を示す．
    \end{flushleft}
\end{sidewaystable}
\clearpage

\subsection{結果の要約}
各手法の特性を以下にまとめる：

\begin{table}[htbp]
    \centering
    \caption{各配車手法の性能評価および総評}
    \label{tab:method_evaluation_summary}
    % resizeboxで表全体をテキスト幅に合わせる
    \resizebox{\textwidth}{!}{%
        \begin{tabular}{lccccl}
            \toprule
            手法 & 重症系RT & 全体RT & 重症6分率 & 全体13分率 & 総評 \\
            \midrule
            直近隊運用 & ○ & ○ & ○ & ○ & 全指標で安定した性能 \\
            傷病度考慮 & ◎ & × & ◎ & × & 重症系に最適だが全体性能は大幅悪化 \\
            MEXCLP & × & × & × & ◎ & 閾値達成率重視，平均RTは悪化 \\
            通常PPO & × & △ & × & △ & 高次元学習の困難さを示す \\
            ハイブリッドPPO & △ & × & △ & × & 期待した効果が得られず \\
            コンパクトPPO & ○ & ◎ & ○ & ○ & 最もバランスの取れた性能 \\
            \bottomrule
        \end{tabular}%
    }
    \footnotesize
    \begin{flushleft}
    注：◎：最良，○：良好，△：中程度，×：不良
    \end{flushleft}
\end{table}

主要な知見：
\begin{itemize}
  \item 優先順位のトレードオフ：傷病度考慮運用の結果に示すとおり，重症系事案への対応を明示的に優先することで，全体のRTと13分閾値達成率が大幅に悪化する．重症系事案は全体の約7\%であるため，残り93\%の軽傷系事案への影響が全体指標に強く反映される
  \item 平均RTと閾値達成率のトレードオフ：MEXCLPの結果が示すとおり，平均RTの悪化と閾値達成率の向上はトレードオフの関係にあることを示している．カバレッジを優先する配車は，平均的な対応速度を犠牲にしつつ，極端に遅い対応を減少させる効果を持つ
  \item 期間における手法間差異：閑散期では各手法の差が小さいが，繁忙期では手法間の差が顕著に拡大する．傷病度考慮運用は配車選択の余地が大きい閑散期で有効性を示した．一方で繁忙期では直近隊運用のような単純で安定した手法の有効性が確認され，直近隊運用に近い運用と考えらえるコンパクトPPOについても同様に有効であった．
  \item 学習の課題：通常PPO及びハイブリッドPPO（999次元）とコンパクトPPO（37次元）の性能差は，高次元状態空間での学習の困難さを示している．次元削減により考慮する情報が削減したことで，学習の方向性が明確化され，実用的な配車性能を獲得したことが考えられる．
\end{itemize}
\clearpage

%%%%%%%%%%%%%%%%%%%%%%%%%  5.6   %%%%%%%%%%%%%%%%%%%%%%%%%%%%%%%%%%%%%%%%%
\section{考察}
\subsection{各手法の特性と適用場面}
本実験の結果から，各配車手法の特性と適切な適用場面について考察する．

直近隊運用は，実装の単純さと安定した性能から，現行の標準的な運用として妥当であることが確認された．すべての評価指標において中程度以上の性能を示し，特に繁忙期における堅牢性が確認された．ただし，重症系事案への優先対応機能を持たないため，重症系RTと6分閾値達成率については改善の余地がある．

傷病度考慮運用は，重症系事案への対応において最も優れた性能を示した．閑散期における重症系RTの0.5～0.7分に及ぶベースラインからの改善は，救命率の向上に寄与する可能性がある．一方，全体RTの大幅な悪化は運用上の課題となる．この手法は，重症系事案への対応を最優先するという政策判断がなされた場合には適用を検討する余地がある．

MEXCLPは，先行研究で報告された特性，すなわち平均RTの悪化と閾値達成率の向上というトレードオフを再現した．本研究の結果から，MEXCLPは平均RTを重視する運用には適さないものの，「一定時間以内の到着保証」を重視する場面では有効な選択肢となりえる．

深層強化学習手法については，コンパクトPPOのみが実用的な性能を達成した．通常PPO及びハイブリッドPPOの結果は，高次元状態空間における学習の困難さ，および報酬設計の複雑さがPPOの性能を制限することを示している．

\subsection{コンパクトPPOの学習効果に関する考察}
コンパクトPPOが直近隊運用と極めて類似したし性能を示したことは，注目に値する．両手法の全体RTの差は最大でも0.19分，13分閾値達成率の差は0.6ポイント以内であり，統計的に優位な差とは言い難い．この結果について，以下の解釈が考えられる．

コンパクトPPOの状態空間には，Top-K救急隊の「移動時間」が主要な特徴量として含まれる．学習の結果，エージェントは移動時間が最も短い救急隊を選択する方策，すなわち「移動時間行列に基づく直近隊運用」を獲得した可能性が高い．

ここで重要な点は，本研究のベースラインである直近隊運用が「移動距離行列」に基づいて最短距離の救急隊を選択しているのに対し，コンパクトPPOは「移動時間行列」に基づく選択を学習した可能性があることである．移動距離と移動時間は高い相関を持つものの，道路状況や経路特性により乖離が生じる場合がある．

実験結果において，コンパクトPPOが直近隊をわずかに上回る期間が存在することは，移動時間ベースの配車が移動距離ベースよりもRTの最小化に適している可能性を示唆する．この解釈が正しければ，深層強化学習が直近隊運用という基本戦略を再発見しつつ，その実装を移動時間ベースに最適化したと解釈できる．

\subsection{深層強化学習アプローチの課題}
本研究における深層強化学習手法の結果は，救急隊配車問題への適用における課題を明らかにした．

第一に，状態空間の設計が学習効率に大きく影響することが確認された．999次元の通常PPO，ハイブリッドPPOは実用的な性能を達成できなかったのに対し，37次元でのコンパクトPPOは直近隊運用と同等以上の性能を示した．救急隊配車問題において，全救急隊の情報を状態として保持することは学習の観点から非効率であり，Top-K救急隊に限定した状態表現が有効であったと考えられる．

第二に，報酬設計の困難さが挙げられる．傷病度考慮運用の有効性に倣い，重症系と軽症系で異なる目標時間を設定し，カバレッジ維持のインセンティブを組み込んだ報酬関数を設計したが，ハイブリッドPPOの結果が示すように複合的な目標を持つ報酬設計は学習を複雑化させる．結果としてカバレッジ維持と，重症系事案への即応の両者を効果的に達成することは実現できなかった．

第三に，学習環境とテスト環境に対する汎化性能の課題がある．通常PPO及びハイブリッドPPOは学習期間に対応する評価期間においては，直近隊運用や傷病度考慮運用を上回る性能を示していた．しかしながら，2024年の多様なテスト期間においてはもちろん，学習期間と近い傾向を持つと考えられる冬季（2月）においても，学習時の性能を再現することはできなかった．これは典型的な過学習による汎化の失敗を示している．本研究では高次元の状態空間に対してハイパーパラメータの調整を試みたものの改善には至らず，コンパクトPPOの結果が示すように，配車判断に必要な特徴量に絞った状態空間において，更なる実験を重ねる必要がある．

\FloatBarrier 
